% Independent continuation by /root/frobenius.
\paragraph{C.4v (Proved: full real multiplication removes every polarization gcd obstruction in characteristic zero).}
Let \(A\) be an abelian variety of dimension \(g\) over an
algebraically closed field of characteristic zero. Suppose that
there is an action
\[
                  \mathcal O_F\hookrightarrow\operatorname{End}(A)
\]
of the full ring of integers of a totally real field \(F\) of
degree \(g\), and a polarization whose Rosati involution fixes
\(F\) pointwise. Then
\[
                              \delta(A)=1.
\]
In particular, if \(g\le3\), \(A\) is an actual direct summand
of a product of Jacobians. This includes a simple abelian
surface or threefold with
\(\operatorname{End}(A)=\mathcal O_K\) for a CM field \(K\)
of degree \(2g\). No principal polarization on \(A\) is assumed.

There is a useful local refinement. Suppose an order
\(\mathcal O\subset F\), rather than the full integer ring,
acts with the same Rosati property, and put
\(c=[\mathcal O_F:\mathcal O]\). Then every prime divisor
of \(\delta(A)\) divides \(c\). Consequently an abelian
surface with \(c\) odd, or an abelian threefold with
\(\gcd(c,6)=1\), is still an actual direct summand of a
product of Jacobians. These conclusions require local
maximality only at two, or at two and three, respectively.

\begin{proof}
We first work over \(\mathbb C\). Put
\(\Lambda=H_1(A,\mathbb Z)\) and let \(J\) denote the complex
structure on \(\Lambda_{\mathbb R}\).
The \(\mathcal O_F\)-module \(\Lambda\) is torsion-free of
rank two, hence projective over this Dedekind domain.
Write \(\mathfrak D_F^{-1}\) for the trace-dual fractional
ideal of \(\mathcal O_F\). The module of integral alternating
forms on \(\Lambda\) satisfying
\[
                         E(ax,y)=E(x,ay)
                  \qquad(a\in\mathcal O_F)
\]
is naturally
\begin{equation}\label{eq:c4-rm-form-module}
 \mathcal P=
 \operatorname{Hom}_{\mathcal O_F}
       \bigl(\textstyle\bigwedge^2_{\mathcal O_F}\Lambda,
                         \mathfrak D_F^{-1}\bigr).
\end{equation}
It is an invertible \(\mathcal O_F\)-module.
Explicitly, an \(F\)-valued alternating form \(\psi\) in
\(\mathcal P\) gives \(E=\operatorname{Tr}_{F/\mathbb Q}\psi\).
Conversely, the nondegenerate trace pairing defines \(\psi\)
by
\[
          \operatorname{Tr}_{F/\mathbb Q}(a\psi(x,y))
                              =E(ax,y)\qquad(a\in F).
\]
The balancing condition makes it \(F\)-bilinear and alternating.
Integrality for every \(a\in\mathcal O_F\) is precisely the
codifferent condition in \eqref{eq:c4-rm-form-module}.

Every real form obtained from \(\mathcal P\) has Hodge type
\((1,1)\). Indeed,
\[
       \Lambda_{\mathbb R}=\bigoplus_{\sigma:F\hookrightarrow\mathbb R}
                                    W_\sigma,\qquad
                         \dim_{\mathbb R}W_\sigma=2.
\]
Balancing makes distinct \(W_\sigma\)'s orthogonal.
The action of \(F\) commutes with \(J\), so each real
two-plane is \(J\)-stable. Since \(J^2=-1\), its determinant
on that plane is one, and every alternating form on the
plane is \(J\)-invariant. This proves the assertion.
Let \(E_0\in\mathcal P\) be the integral positive Riemann
form of the given polarization; its balancing condition
is exactly the stated Rosati hypothesis.

Fix a rational prime \(\ell\), including \(\ell=2\).
Over the product of discrete valuation rings
\(\mathcal O_F\otimes\mathbb Z_\ell\), the module
\(\Lambda_\ell\) is free of rank two and
\(\mathcal P_\ell\) is free of rank one. Choose a generator
\(\psi_\ell\) of the latter.
In a basis of \(\Lambda_\ell\), its trace form is
\[
 ((x_1,x_2),(y_1,y_2))
       \longmapsto
 \operatorname{Tr}_{F\otimes\mathbb Q_\ell/\mathbb Q_\ell}
                    \bigl(b(x_1y_2-x_2y_1)\bigr),
\]
where \(b\) generates the local codifferent. The definition
of the codifferent makes this pairing perfect over
\(\mathbb Z_\ell\).
Choose \(E_1\in\mathcal P\) whose reduction modulo \(\ell\)
agrees with this form. This is possible because
\(\mathcal P/\ell\mathcal P
 \simeq\mathcal P_\ell/\ell\mathcal P_\ell\).
For a sufficiently large positive integer \(n\),
\[
                              E=E_1+\ell n E_0
\]
is positive definite in the Riemann-form sense, while its
reduction modulo \(\ell\) is still perfect. Positivity here
is an elementary fact about a fixed real symmetric form
plus a sufficiently large multiple of a positive definite
one; no local-to-global principal-polarization assertion
is being made.

Appell--Humbert and the positivity criterion produce an
ample algebraic line bundle \(L_\ell\) with this integral
Riemann form. The exact primary statement used is the
proof of \citet[Proposition 6.1, pp.~46--47]{frC4MilneCM2020RM},
where both the Appell--Humbert theorem and the Lefschetz
positivity theorem are stated. The induced homomorphism
on integral homology is the map given by \(E\), so
\(\deg\varphi_{L_\ell}=|\det E|\) is prime to \(\ell\);
see also the displayed polarization map on p.~26 of
the same source. Since
\(\deg\varphi_{L_\ell}=\chi(L_\ell)^2\), we have
\(\ell\nmid\chi(L_\ell)\). We construct this line bundle
separately for each prime. Hence no prime divides
\(\delta(A)\), proving \(\delta(A)=1\).

For an arbitrary algebraically closed characteristic-zero
field, define \(A\), the finite set of endomorphisms giving
the action, and the polarization over a finitely generated
subfield, and embed it into \(\mathbb C\).
For each \(\ell\), the resulting ample line bundle spreads
over a nonempty scheme of finite type over that subfield.
After shrinking, ampleness and its Euler characteristic
are preserved. A closed point has a finite residue extension,
which embeds into the given algebraically closed field.
Its line bundle therefore has the same Euler characteristic
prime to \(\ell\). This proves the same gcd assertion.
C.4j then constructs the actual retraction for \(g\le3\).

Finally, when \(\operatorname{End}(A)=\mathcal O_K\) for
the indicated CM field, take \(F=K^+\).
Any Rosati involution on \(K\) is complex conjugation:
it is positive, and
\citet[Proposition 1.39 and Corollary 1.40,
pp.~20--21]{frC4MilneCM2020RM}
identify the positive involution on a CM field uniquely.
Its restriction to \(F\) is the identity, and
\(\mathcal O_F\subset\mathcal O_K\). Thus all the stated
hypotheses hold.

For the local refinement, let \(\mathcal O\) be the given
order and replace \(\mathcal P\) by the integral lattice
\[
 \mathcal P_{\mathcal O}=
 \{E\in\operatorname{Hom}_{\mathbb Z}
       (\textstyle\bigwedge^2_{\mathbb Z}\Lambda,\mathbb Z):
             E(ax,y)=E(x,ay)\text{ for all }a\in\mathcal O\}.
\]
Using a finite integral basis of the order, these are
finitely many integral linear equations. Localization
is exact, so this lattice tensored with \(\mathbb Z_\ell\)
is exactly the lattice of locally balanced integral
forms. If \(\ell\nmid c\), then
\(\mathcal O\otimes\mathbb Z_\ell=
\mathcal O_F\otimes\mathbb Z_\ell\). The local homology
module is therefore free of rank two over each local
integer ring, and the codifferent calculation above
again supplies a perfect local form.
Lift its residue modulo \(\ell\) to
\(\mathcal P_{\mathcal O}\) and add \(\ell nE_0\).
The real-plane and positivity arguments are unchanged,
since \(\mathcal O\otimes\mathbb Q=F\).
We obtain an ample line bundle whose Euler
characteristic is prime to \(\ell\). Thus
\[
            \operatorname{rad}\delta(A)
                   \mid\operatorname{rad}c.
\]
For surfaces and threefolds we combine this with the
already proved universal bounds \(s_J(A)\mid2\) and
\(s_J(A)\mid6\), respectively, and with
\(s_J(A)\mid\delta(A)\). This proves the stated
local refinements; it does not assert that
\(\delta(A)=1\) at a prime where the order is nonmaximal.

In the simple CM case the same refinement applies if
the endomorphism order in \(K\) is locally maximal at
the indicated obstruction primes. Indeed its
intersection with \(K^+\) is then locally
\(\mathcal O_{K^+}\), and the positive Rosati
involution still restricts to the identity there.
\end{proof}

The same affirmative conclusion holds for any good
reduction of one of these varieties of dimension at most
three. Indeed, \(\delta(A)=1\) supplies finitely many
ample line bundles whose Euler characteristics have gcd
one. After a finite extension of a complete discrete
valuation ring of good reduction, all these line bundles
extend to the smooth, hence regular, abelian scheme by Cartier
closure. Its specialization need not stay ample, but
its Euler characteristic stays constant in this proper
flat family. The definition of \(\delta\) using all line
bundles therefore still gives \(\delta=1\) on the
geometric special fibre; the ample-gcd equivalence in
C.4j supplies the required retraction in every residue
characteristic. We retain the full-order and Rosati
hypotheses for the assertion \(\delta=1\).
The locally maximal refinement also specializes:
the finitely many line bundles needed at the
obstruction primes extend and preserve their Euler
characteristics by the same argument. We make no
assertion at a nonmaximal obstruction prime, nor
for the generic surface with endomorphism ring
\(\mathbb Z\).
